\documentclass[11pt]{article}
\usepackage[margin=1in]{geometry}
\usepackage{amsmath,amssymb,amsthm,bm,mathtools}
\usepackage{algorithm}
\usepackage{algpseudocode}
\usepackage{booktabs}
\usepackage{enumitem}
\usepackage{hyperref}
\hypersetup{colorlinks=true,linkcolor=blue,citecolor=blue,urlcolor=blue}

\title{Fairness-Aware Regression via Composite \\Nonsmooth Optimization}
\author{Draft}
\date{\today}

% --- Theorem environments ---
\theoremstyle{plain}
\newtheorem{theorem}{Theorem}
\newtheorem{proposition}{Proposition}
\newtheorem{lemma}{Lemma}
\theoremstyle{definition}
\newtheorem{definition}{Definition}
\newtheorem{assumption}{Assumption}
\newtheorem{remark}{Remark}

% --- Macros ---
\newcommand{\R}{\mathbb{R}}
\newcommand{\E}{\mathbb{E}}
\newcommand{\prox}{\mathrm{prox}}
\newcommand{\dom}{\mathrm{dom}}
\newcommand{\argmin}{\mathop{\mathrm{arg\,min}}}
\newcommand{\argmax}{\mathop{\mathrm{arg\,max}}}
\newcommand{\ip}[2]{\left\langle #1,#2 \right\rangle}
\newcommand{\norm}[1]{\left\lVert #1 \right\rVert}
\newcommand{\1}{\mathbf{1}}
\newcommand{\ind}{\mathbb{I}}
% convex-analysis indicator (0 on set, +inf outside)
\newcommand{\indC}[1]{\iota_{#1}}

\begin{document}
\maketitle

\begin{abstract}
We develop a unified framework for \emph{fairness-aware regression} in which a smooth data-fidelity term is combined with a possibly nonsmooth or even discontinuous fairness functional. We: (i) motivate fairness in regression from predictive parity and allocative harm perspectives; (ii) define a general fairness operator that subsumes mean (prediction or residual) parity, max-deviation gaps, max--min residual gaps, absolute-error parity, variance/quantile parity, and indicator-type constraints; (iii) formalize the optimization template
\[
\psi(\lambda):=\min_{t\in\R^p}\; g(t)+\lambda\,h(c(t)),
\]
with $g$ differentiable and $h$ proper closed (nonsmooth allowed), $c$ a smooth map from parameters to fairness-relevant statistics; (iv) derive proximal and primal--dual algorithms, including prox-gradient (Case A), prox-linear with Chambolle--Pock inner loops (Case B), and ADMM for constrained fairness; (v) provide value-function identities and general sensitivity bounds that quantify the accuracy--fairness trade-off, including Lipschitz and curvature-type estimates under mild regularity; and (vi) add a \emph{proportional-fair} (Nash social welfare) regression that stays convex by operating on error headroom. Detailed pseudocode and RMSE/MSE examples are provided.
\end{abstract}

\section{Motivation: Why Fairness in Regression?}
Classic fairness literature often focuses on classification; yet numerous decisions are driven by \emph{continuous} predictions: credit limits, insurance premiums, tax assessments, energy rates, health risk scores, and salary recommendations. Disparities in \emph{residuals} (systematic under- or over-prediction by group) are particularly harmful:
\begin{itemize}[leftmargin=2em]
\item \textbf{Allocative harm:} Persistent overestimation of property values for a subpopulation implies regressive taxation; persistent underestimation of health risks for a region implies under-allocation of resources.
\item \textbf{Procedural equity:} Even if aggregate error is low, skewed residuals erode trust and legitimacy in public-facing models.
\item \textbf{Robustness:} Penalizing inter-group disparities reduces sensitivity to dataset shifts that differentially affect subpopulations.
\end{itemize}
Thus, for regression we often target \emph{parity of prediction means}, \emph{parity of residual means}, \emph{parity of absolute residuals} (MAE), or constraints on \emph{max/min residuals}, among others.

\section{General Fairness Functional}
Let $\{(x_i,y_i,g_i)\}_{i=1}^n$ with group labels $g_i\in\{1,\dots,G\}$, and parametric regressor $\hat y_i(t)=f(x_i;t)$.

We define a general fairness operator as a composition:
\[
 c(t) \;=\; \Gamma\!\left( \Phi\big(\hat y(t),y\big) \right) \in \R^m,
\]
where:
\begin{itemize}[leftmargin=2em]
 \item $\Phi:\R^n\times\R^n\to\R^q$ transforms predictions and labels into per-sample or per-group \emph{statistics} (e.g., residuals $r=\hat y-y$, absolute residuals $|r|$, squared residuals $r^2$, or indicators $\ind\{|r|>\tau\}$).
 \item $\Gamma:\R^q\to\R^m$ aggregates by groups and optionally differences groups (e.g., $S$ is a group-averaging operator; $K$ encodes pairwise differences).
\end{itemize}
The fairness scalarization $h:\R^m\to(-\infty,+\infty]$ is then applied to $c(t)$. We assume henceforth $h(u)\ge 0$ and $h(0)=0$ (true for norms and many indicators), which ensures certain monotonicity properties later.

\paragraph{Examples (all fit $h(c(t))$):}
\begin{enumerate}[leftmargin=2em,label=\textbf{E\arabic*.}]
 \item \textbf{Pairwise mean \emph{prediction} gap:}
 $c(t)=K S \hat y(t)$ and $h(u)=\|u\|_\infty$ or $\|u\|_1$.

 \item \textbf{Pairwise mean \emph{residual} gap (systematic bias):}
 $c(t)=K S(\hat y(t)-y)$ and $h(u)=\|u\|_\infty$ (max absolute gap) or $\|u\|_1$.

 \item \textbf{Max--min residual control:}
 Let $r_i(t)=\hat y_i(t)-y_i$. Define
 \[
 c(t)=\begin{bmatrix} \max_i r_i(t) - \min_i r_i(t) \end{bmatrix},
 \quad
 h(u)=u \quad(\text{penalty}) \ \ \text{or}\ \ h(u)=\indC{\{u\le \tau\}}(u).
 \]
 A convex surrogate: $c(t) = \big\| r(t)-\bar r \,\1 \big\|_{\infty}$ (deviation from mean).

 \item \textbf{Mean absolute error (MAE) parity:}
 $c(t)=K S|r(t)|$ with $h(u)=\|u\|_\infty$ or $\|u\|_1$.

 \item \textbf{Variance/dispersion parity:}
 $c(t)=K S\big((r(t)-\overline r_g)^2\big)$ (group-wise residual variance differences). This is nonconvex in $t$ because of the $-\overline r_g^2$ term; we typically use weak convexity or smooth surrogates.

 \item \textbf{Quantile parity:}
 $c(t)=K\,\big(Q_\alpha(r\,|\,g)\big)_{g=1}^G$ with $Q_\alpha$ a (Huber-smoothed) quantile operator; $h$ as $\ell_\infty$ or $\ell_1$.

 \item \textbf{Hard constraints:} $h(c)=\indC{\{\|c\|_\infty\le\tau\}}(c)$ or $\indC{\{c\le \tau\}}(c)$, producing discontinuities; handled via projections/ADMM.
\end{enumerate}

\section{Optimization Framework}
We study
\begin{equation}\label{eq:master}
 \min_{t\in\R^p}\; F_\lambda(t) \;=\; g(t) \;+\; \lambda\, h\big(c(t)\big),
\end{equation}
with:
\begin{itemize}[leftmargin=2em]
 \item $g$ \emph{differentiable}, $\nabla g$ $L$-Lipschitz. Typical: MSE or RMSE surrogate
 \[
 g_{\text{MSE}}(t)=\frac{1}{2n}\sum_{i=1}^n\big(\hat y_i(t)-y_i\big)^2,\quad
 \nabla g_{\text{MSE}}(t)=\frac{1}{n}\sum_i \big(\hat y_i(t)-y_i\big)\,\nabla_t f(x_i;t).
 \]
 (RMSE is concave in squared errors; minimizing it is nonconvex. We therefore optimize MSE.)
 \item $c$ \emph{smooth}: composition of model $f(\cdot;t)$ with linear operators ($S$, $K$) and possibly elementwise maps (absolute value admits weak convexity; Huber smoothing is optional).
 \item $h$ \emph{proper closed}, possibly nonsmooth or discontinuous (indicators). Often convex (e.g., norms, max, support functions).
\end{itemize}

Two regimes:
\begin{description}[leftmargin=2em,style=nextline]
 \item[Case A:] $c(t)=t$ so $F_\lambda(t)=g(t)+\lambda h(t)$ (direct composite).
 \item[Case B:] $c$ nontrivial (output-composite); $h\circ c$ has no simple prox.
\end{description}

\section{Algorithms}
\subsection{Case A: Proximal Gradient and Acceleration}
Assume $h$ convex. With step $\alpha\in(0,1/L]$,
\[
 t^{k+1}=\prox_{\alpha\lambda h}\!\big(t^k-\alpha\nabla g(t^k)\big).
\]
ISTA yields $O(1/k)$ decrease (convex); FISTA (Nesterov) achieves $O(1/k^2)$. If $g$ is $\mu$-strongly convex, rates are linear. For weakly nonconvex $h$ (e.g., capped-$\ell_1$), the proximal gradient mapping $\tfrac{1}{\alpha}\big(t-\prox_{\alpha\lambda h}(t-\alpha\nabla g(t))\big)$ converges to $0$.

\paragraph{Prox toolbox (examples).}
\begin{itemize}[leftmargin=2em]
 \item $\ell_1$: soft-thresholding; $\prox_{\alpha\lambda \|\cdot\|_1}(z)=\mathrm{sign}(z)\odot\max\{|z|-\alpha\lambda,0\}$.
 \item $\ell_2^2$: shrinkage; $\prox_{\alpha(\lambda/2)\|\cdot\|_2^2}(z)=\frac{1}{1+\alpha\lambda}z$.
 \item Group $\ell_{2,1}$: block soft-thresholding on groups.
 \item Indicator of convex set $\mathcal{C}$: projection $\Pi_{\mathcal{C}}$.
\end{itemize}

\subsection{Case B: Prox-Linear Outer Loop + Primal--Dual Inner Loop}
We linearize $c$ at $t^k$:
\[
 c(t)\approx c(t^k)+J_c(t^k)(t-t^k).
\]
Let $q^k=t^k-\alpha\nabla g(t^k)$, $A^k=J_c(t^k)$, $b^k=c(t^k)-A^k t^k$.
The \emph{prox-linear subproblem} is
\begin{equation}\label{eq:proxlin}
 \min_{t}\; \frac{1}{2\alpha}\norm{t-q^k}^2 + \lambda\,h\!\big(A^k t + b^k\big).
\end{equation}
It is a convex composite with a linear map into $h$ and can be solved by:
\begin{itemize}[leftmargin=2em]
 \item \textbf{Chambolle--Pock (CP)}: primal--dual with dual variable $y$ for $h$,
 \begin{align*}
 y^{\ell+1}&=\prox_{\sigma (\lambda h)^*}\big( y^\ell + \sigma(A^k \bar t^\ell + b^k)\big),\\
 t^{\ell+1}&=\frac{1}{1+\tau/\alpha}\Big( t^\ell - \tau (A^k)^\top y^{\ell+1} + (\tau/\alpha) q^k \Big),\\
 \bar t^{\ell+1}&=t^{\ell+1}+\theta(t^{\ell+1}-t^\ell),
 \end{align*}
 with $\tau\sigma\|A^k\|^2<1$.
 \item \textbf{ADMM form}: introduce $z=A^k t+b^k$ and separate $t$ and $z$ updates (the $z$-update uses $\prox_{\lambda h}$).
\end{itemize}
Under standard conditions, prox-linear converges globally to critical points; in convex settings one gets sublinear rates for stationarity measures, and linear rates under error bounds or strong convexity in $g$.

\subsection{Constrained Fairness via ADMM}
For $\min g(t)$ s.t. $h(c(t))\le \tau$, write
\[
 \min_{t,z}\; g(t)+\indC{\{h(z)\le \tau\}}(z)\quad\text{s.t.}\quad z=c(t).
\]
The augmented Lagrangian with scaled dual $u$ and penalty $\rho>0$ leads to iterations:
\begin{align*}
 t^{k+1}&\approx \argmin_t \; g(t)+\frac{\rho}{2}\norm{c(t)-z^k+u^k}^2,\\
 z^{k+1}&=\Pi_{\{h\le\tau\}}\big(c(t^{k+1})+u^k\big),\\
 u^{k+1}&=u^k + c(t^{k+1})-z^{k+1}.
\end{align*}
The $t$-update uses (stochastic) gradients and JVP/VJP for $c$.

\section{Fairness Operators Instantiated}
We now detail common $c$ and $h$ pairs and their conjugate/prox.

\subsection{Pairwise Mean \emph{Residual} Gap (max absolute gap)}
Let $r(t)=\hat y(t)-y\in\R^n$, $S\in\R^{G\times n}$ group-averages, $K\in\R^{m\times G}$ pairwise differences. Then
\[
 c(t)=K S r(t), \qquad h(u)=\|u\|_\infty.
\]
Conjugate $h^*(y)=\indC{\{\|y\|_1\le 1\}}(y)$. Hence
\[
 \prox_{\sigma (\lambda h)^*}(w)=\Pi_{\{\|y\|_1\le \lambda\}}(w),
\]
projection onto the $\ell_1$ ball (sort-threshold).

\subsection{Mean Absolute Error Parity}
$c(t)=K S|r(t)|$, $h(u)=\|u\|_\infty$. The absolute value makes $c$ piecewise-smooth and $h\circ c$ weakly convex; prox-linear is still applicable; Huber smoothing of $|r|$ is an alternative.

\subsection{Max--Min Residual Gap}
$c(t)=[\max_i r_i(t)-\min_i r_i(t)]$. Convex surrogate: $c(t)=\|r(t)-\overline r \,\1\|_\infty$. With $h(u)=u$ or $h(u)=\indC{\{u\le \tau\}}$, we maintain convexity.

\section{Proportional-Fair (Nash) Regression for Continuous Outcomes}
Proportional fairness (PF) maximizes $\sum_i \log u_i$ over a convex set. For regression, choose per-sample \emph{error headroom} $u_i=\tau - t_i$ with $t_i\ge |r_i|$. Then the \emph{individual PF} program is
\begin{align}
 \max_{\beta,b,\,t\in\R^n}\; &\sum_{i=1}^n \log(\tau - t_i) \label{eq:pf}\quad\\
 \text{s.t. }\; & t_i \ge y_i - x_i^\top\beta - b,\quad t_i \ge -y_i + x_i^\top\beta + b,\quad 0\le t_i \le \tau-\epsilon,\; \forall i,\\
 & \|\beta\|_2 \le B \quad (\text{optional size control}).\nonumber
\end{align}
\begin{proposition}[Convexity and existence]
Problem~\eqref{eq:pf} is a concave maximization over a convex feasible set (equivalently, minimize a convex function by negating the objective); hence any local optimum is global. If $\tau>0$ and the residual constraints are feasible, a solution exists. The KKT multipliers on $t_i$ quantify \emph{proportional} trade-offs.
\end{proposition}
A \emph{group PF} variant replaces $\log(\tau - t_i)$ by $\sum_g w_g \log(\tau_g - e_g)$ with $e_g$ the mean absolute error in group $g$.

\paragraph{Log-target consistency and ratios.} If training is conducted in log space, constraining $|\log y_i - \eta_i|\le \rho$ with $\eta_i=x_i^\top\beta$ enforces multiplicative bounds $e^{-\rho}\le \hat y_i/y_i\le e^{\rho}$ after exponentiation; PF can be applied to the resulting headrooms $\tau-\!|\log y_i-\eta_i|$.

\section{Pseudocode}
\subsection{ISTA / FISTA (Case A)}
\begin{algorithm}[H]
\caption{ISTA for $F_\lambda(t)=g(t)+\lambda h(t)$}
\begin{algorithmic}[1]
\State \textbf{Input:} $t^0$, step $\alpha\in(0,1/L]$ or backtracking, tol $\varepsilon$
\For{$k=0,1,2,\dots$}
  \State $z \gets t^k - \alpha \nabla g(t^k)$
  \State $t^{k+1} \gets \prox_{\alpha\lambda h}(z)$
  \If{$\frac{1}{\alpha}\norm{t^{k+1}-t^k}\le \varepsilon$} \textbf{break} \EndIf
\EndFor
\end{algorithmic}
\end{algorithm}

\begin{algorithm}[H]
\caption{FISTA (monotone, with restart) for $g+\lambda h$}
\begin{algorithmic}[1]
\State \textbf{Input:} $t^0$, $y^0=t^0$, $s_0=1$, $\alpha$, $\varepsilon$
\For{$k=0,1,2,\dots$}
  \State $z \gets y^k - \alpha \nabla g(y^k)$
  \State $t^{k+1} \gets \prox_{\alpha\lambda h}(z)$
  \State $s_{k+1}\gets \frac{1+\sqrt{1+4 s_k^2}}{2}$,\quad
  $y^{k+1}\gets t^{k+1} + \frac{s_k-1}{s_{k+1}}(t^{k+1}-t^k)$
  \If{$F_\lambda(t^{k+1})>F_\lambda(t^k)$ \textbf{or} $\ip{t^{k+1}-t^k}{t^k-t^{k-1}}>0$}
    \State $y^{k+1}\gets t^{k+1}$,\quad $s_{k+1}\gets 1$ \Comment{restart}
  \EndIf
  \If{$\frac{1}{\alpha}\norm{t^{k+1}-t^k}\le \varepsilon$} \textbf{break} \EndIf
\EndFor
\end{algorithmic}
\end{algorithm}

\subsection{Prox-Linear + Chambolle--Pock (Case B)}
\begin{algorithm}[H]
\caption{Prox-Linear Outer with CP Inner for $g(t)+\lambda h(c(t))$}
\begin{algorithmic}[1]
\State \textbf{Input:} $t^0$, outer step $\alpha>0$, tolerances $\varepsilon_{\rm out},\varepsilon_{\rm in}$
\For{$k=0,1,2,\dots$}
  \State $q \gets t^k - \alpha \nabla g(t^k)$
  \State Compute $A \gets J_c(t^k)$ (through JVP/VJP), $b\gets c(t^k)-A t^k$
  \State \textbf{Inner:} Solve $\min_t \frac{1}{2\alpha}\norm{t-q}^2 + \lambda h(A t+b)$ by CP
  \State \quad Choose $\tau,\sigma>0$ with $\tau\sigma\|A\|^2<1$, set $(t^0_{\rm in},y^0)=(t^k,0)$
  \For{$\ell=0,1,2,\dots$}
     \State $y^{\ell+1} \gets \prox_{\sigma(\lambda h)^*}(y^\ell + \sigma (A \bar t^\ell + b))$
     \State $t^{\ell+1} \gets \frac{1}{1+\tau/\alpha}\Big(t^\ell - \tau A^\top y^{\ell+1} + (\tau/\alpha) q \Big)$
     \State $\bar t^{\ell+1}\gets t^{\ell+1} + \theta (t^{\ell+1}-t^\ell)$
     \If{primal--dual gap $\le \varepsilon_{\rm in}$} \textbf{break} \EndIf
  \EndFor
  \State $t^{k+1}\gets t^{\ell+1}$
  \If{$\frac{1}{\alpha}\norm{t^{k+1}-t^k}\le \varepsilon_{\rm out}$} \textbf{break} \EndIf
\EndFor
\end{algorithmic}
\end{algorithm}

\subsection{Constrained Fairness via ADMM}
\begin{algorithm}[H]
\caption{ADMM for $\min g(t)$ s.t. $h(c(t))\le \tau$}
\begin{algorithmic}[1]
\State \textbf{Input:} $t^0,z^0,u^0$, penalty $\rho>0$
\For{$k=0,1,2,\dots$}
  \State $t^{k+1}\gets \argmin_t \; g(t)+\frac{\rho}{2}\norm{c(t)-z^k+u^k}^2$ \Comment{gradients + JVP/VJP}
  \State $z^{k+1}\gets \Pi_{\{h\le \tau\}}\big(c(t^{k+1})+u^k\big)$
  \State $u^{k+1}\gets u^k + c(t^{k+1})-z^{k+1}$
  \If{primal/dual residuals below tol} \textbf{break} \EndIf
\EndFor
\end{algorithmic}
\end{algorithm}

\section{Constrained vs. Penalized: Slope Law and Selection Rules}
Define value functions
\[
 \psi(\lambda)=\inf_t \{ g(t)+\lambda h(c(t))\},\qquad
 \phi(\tau)=\inf_t \{ g(t): h(c(t))\le \tau\}.
\]
Assume convexity (of $g$ and $h\circ c$) and Slater’s condition.

\begin{proposition}[Envelope identities \,\&\, monotonicity]
Let $t_\lambda\in\argmin \{g(t)+\lambda h(c(t))\}$ and $t_\tau$ be optimal for the constrained form with multiplier $\lambda^*_\tau\ge 0$. If $h\ge 0$ and $h(0)=0$, then
\begin{align}
 &\psi \text{ is convex and nondecreasing in } \lambda,\quad h(c(t_\lambda))\in \partial \psi(\lambda)\; (\psi'(\lambda)=h(c(t_\lambda)) \text{ when single-valued}),\\
 &\phi \text{ is convex and nonincreasing in } \tau,\quad -\lambda^*_\tau \in \partial \phi(\tau)\; (\phi'(\tau)=-\lambda^*_\tau \text{ when differentiable}).
\end{align}
Moreover $\tau(\lambda):=h(c(t_\lambda))$ is nonincreasing in $\lambda$.
\end{proposition}

\paragraph{Percent-slope identity.} On the frontier, the marginal percent tradeoff obeys
\[
 \frac{d\log g}{d\log h} = -\lambda\, \frac{h}{g},
\]
useful for selecting $\lambda$ to match an allowable "+\%$f$ per \(-\%$g$)" rate.

\subsection{Lipschitz Sensitivity of the Path}
\begin{assumption}\label{ass:reg}
$g$ is $\mu$-strongly convex with $L$-Lipschitz gradient; $c$ is $C^1$ with $\sup\|J_c(t)\|\le M$ on the relevant level set; $h$ is $G$-Lipschitz.
\end{assumption}

\begin{proposition}[Solution and fairness Lipschitzness]\label{prop:lips}
Under Assumption~\ref{ass:reg} and convexity, for penalized solutions $t_\lambda,t_{\lambda'}$:
\[
 \norm{t_\lambda-t_{\lambda'}} \le \frac{G\,M}{\mu}\,|\lambda-\lambda'|,
 \qquad
 \big|h(c(t_\lambda)) - h(c(t_{\lambda'}))\big| \le \frac{G^2 M}{\mu}\,|\lambda-\lambda'|.
\]
Furthermore, by $L$-smoothness,
\[
 |\,g(t_\lambda)-g(t_{\lambda'})\,|
 \le \frac{L}{2}\norm{t_\lambda-t_{\lambda'}}^2 + \norm{\nabla g(t_{\lambda'})}\,\norm{t_\lambda-t_{\lambda'}}.
\]
\end{proposition}

\subsection{Curvature Bounds and Slope Control}
When $h\circ c$ is convex and $g$ is twice differentiable, $\psi$ is convex and a.e. twice differentiable with
\[
 \psi'(\lambda)=\tau(\lambda),\qquad \psi''(\lambda)\ge 0.
\]
If in addition $\lambda\mapsto t_\lambda$ is single-valued and differentiable (e.g., under strong convexity of $g$ and strict complementarity), implicit differentiation yields
\[
 \frac{d \tau}{d\lambda} 
 = \ip{\nabla (h\!\circ\! c)(t_\lambda)}{ \frac{dt_\lambda}{d\lambda} }.
\]
Bounding $\|dt_\lambda/d\lambda\|$ via strong monotonicity and $\|\nabla (h\!\circ\! c)\|\le GM$ gives
\[
 0 \;\le\; -\frac{d\tau}{d\lambda} \;\le\; \frac{G^2 M^2}{\mu}.
\]
Thus fairness decreases smoothly with $\lambda$ at a controlled rate; the frontier is \emph{Lipschitz--smooth}.

\subsection{Certified Slopes from Solvers}
\begin{itemize}[leftmargin=2em]
 \item \textbf{Penalized:} $\tau(\lambda)=h(c(t_\lambda))$ is a subgradient of $\psi$; finite differences yield $\Delta g$ and $\Delta \tau$.
 \item \textbf{Constrained (primal--dual):} The dual variable converges to $\lambda^*_\tau$, which equals $-d\phi/d\tau$; primal--dual \emph{gap} certificates bound the slope interval.
\end{itemize}

\section{RMSE/MSE Examples (Gradients and Jacobians)}
Let $g(t)=\frac{1}{2n}\sum_i (\hat y_i(t)-y_i)^2$. Then
\[
 \nabla g(t)=\frac{1}{n}\sum_i (\hat y_i(t)-y_i)\,\nabla_t f(x_i;t).
\]
For $c(t)=K S (\hat y(t)-y)$, the Jacobian acts as
\[
 J_c(t)\,v = K S \big(J_{\hat y}(t)\,v\big),\qquad
 J_c(t)^\top w = J_{\hat y}(t)^\top (S^\top K^\top w),
\]
which are efficiently computed by JVP/VJP in modern autodiff libraries.

\section{Implementation Blueprint}
\begin{enumerate}[leftmargin=2em]
 \item \textbf{Operators:} Precompute $S$ (group means) and $K$ (pairwise differences). For residual-based parity, use $r=\hat y-y$.
 \item \textbf{Autograd:} Implement JVP/VJP for $t\mapsto \hat y(t)$; never form $J_{\hat y}$ explicitly.
 \item \textbf{Outer step \& backtracking:} Start with $\alpha=1/L$ (backtrack on $g$’s quadratic majorization); ensure decrease of the majorized model in prox-linear.
 \item \textbf{Inner steps:} For CP, choose $\tau,\sigma$ with $\tau\sigma\|A\|^2<1$; estimate $\|A\|$ by 3--5 power iterations (JVP/VJP).
 \item \textbf{Stopping:} Outer stationarity via prox-gradient mapping; inner primal--dual gap or KKT residual. Record $(g,h)$ pairs and dual multipliers.
 \item \textbf{Path tracing:} Either grid in $\lambda$ (warm-started) or bisection to hit target fairness $\tau$ (monotone $F(\lambda)=h(c(t_\lambda))$).
 \item \textbf{Stability:} Add mild $\ell_2$ ridge to $g$ to improve $\mu$ and tighten bounds in Prop.~\ref{prop:lips}.
\end{enumerate}

\section{Numerical Protocol (Suggested)}
\begin{itemize}[leftmargin=2em]
 \item \textbf{Splits:} Stratify train/val/test by groups; report per-group and overall MSE/MAE, plus fairness metrics.
 \item \textbf{Curves:} Plot $g^\star$ vs $\tau$ and $\tau$ vs $\lambda$ with error bars from duality gaps; overlay multipliers (slope) and finite-difference slopes.
 \item \textbf{Baselines:} ERM; penalized with multiple $h$ choices; constrained with multiple $\tau$; smoothed surrogates (Huber/log-sum-exp); proportional-fair.
\end{itemize}

\section{Limitations and Extensions}
\begin{itemize}[leftmargin=2em]
 \item \textbf{Nonconvex models:} Neural nets make $F_\lambda$ nonconvex. Prox-linear still converges to critical points; KL/error-bound properties often induce fast local rates.
 \item \textbf{Small groups \& shift:} Regularize small $|G_g|$ (weights in $S$), consider distributionally robust $g$ (e.g., DRO with Wasserstein ball) to improve transfer.
 \item \textbf{Individual fairness:} Add Lipschitz constraints $|\hat y_i-\hat y_j|\le L\,d(x_i,x_j)$ over a graph of similar pairs.
 \item \textbf{Multi-criteria fairness:} Vector fairness $h(c)=\|(h_1(c),\dots,h_r(c))\|_{\mathcal{N}}$ via norms or multiobjective scalarization is straightforward.
\end{itemize}

\section*{Conclusion}
We provided a complete optimization-first treatment of fairness-aware regression: a general fairness operator $c(t)$, nonsmooth scalarizations $h$, composite optimization algorithms (prox-gradient, prox-linear + primal--dual, ADMM), a convex proportional-fair regression, and rigorous trade-off analysis via value functions, envelope identities, and Lipschitz/curvature bounds. The framework is implementable at scale via JVP/VJP and yields certified slopes along the accuracy--fairness frontier.

\appendix
\section{Appendix A: Proof Sketches}
\paragraph{Envelope identities.} Apply Danskin’s theorem to $\psi(\lambda)=\inf_t \{g(t)+\lambda h(c(t))\}$. For $\phi(\tau)=\inf_t \{g(t)+\indC{\{h(c(t))\le \tau\}}(t)\}$, Lagrangian duality yields $-\lambda^*_\tau\in\partial\phi(\tau)$.
\paragraph{Path Lipschitzness.} Optimality: $\nabla g(t_\lambda)+\lambda\,\xi_\lambda=0$ with $\xi_\lambda\in\partial(h\circ c)(t_\lambda)$. Subtract at $\lambda'$, inner-product with $t_\lambda-t_{\lambda'}$, use strong monotonicity of $\nabla g$ and $\|\xi\|\le GM$ (chain rule) to bound $\|t_\lambda-t_{\lambda'}\|$; then bound $|h(c(\cdot))|$ by $G$-Lipschitzness and $\|J_c\|\le M$.

\section{Appendix B: Projections and Prox Operators}
\begin{itemize}[leftmargin=2em]
 \item \textbf{$\ell_1$-ball projection} (for $\prox_{\sigma (\lambda\|\cdot\|_\infty)^*}$): sort $|w|$, find threshold $\theta$ with $\sum_i \max\{|w_i|-\theta,0\}=\lambda$, then project $w$ to $w\cdot \min\{1,\lambda/\|w\|_1\}$ after thresholding.
 \item \textbf{$\ell_\infty$-box projection} (for hard constraints): clip components to $[-\tau,\tau]$.
 \item \textbf{Group block soft-thresholding} for $\ell_{2,1}$ penalties.
\end{itemize}

\section{Appendix C: Smoothing Alternatives}
For nonconvex or nondifferentiable $\Phi$ (e.g., $|r|$, quantiles), consider Huberized absolute values, softmax/log-sum-exp for max, or Moreau envelope of $h$.
Smoothing allows plain gradient methods but introduces bias; prox-linear keeps the original nonsmooth geometry.

\end{document}
